\chapter{Aspectos Conceituais} 

Para solucionar o problema da extração manual e garantir a confiabilidade dos dados ingeridos pela Veeries, fez-se necessário o uso de padrões arquiteturais consolidados na área de engenharia de dados. Este capítulo apresenta, de forma prática, os conceitos de \textit{pipelines} de dados, a Arquitetura \textit{Medallion} e a utilização de \textit{schemas} como contratos de dados, fundamentando as escolhas técnicas implementadas na solução.

\section{Pipelines e Qualidade de Dados}

Um \textit{pipeline} de dados consiste em um conjunto automatizado de processos que realiza a extração, o processamento, a validação e o carregamento de informações de um sistema de origem para um ambiente de destino \cite{polisetty23}. Em oposição aos processos manuais e rudimentares, a adoção de fluxos automatizados garante a continuidade da informação e mitiga drasticamente a ocorrência de falhas humanas \cite{reis22}.

A qualidade dos dados (\textit{Data Quality}) é o objetivo central dessas rotinas estruturadas, sendo uma medida avaliada com base em fatores como precisão, completude, consistência e confiabilidade \cite{polisetty23}. No contexto do processamento de dados não estruturados e semi-estruturados — como a leitura e extração de tabelas em arquivos PDF —, a automação do fluxo é o primeiro passo obrigatório para estabelecer a confiabilidade do que será entregue à equipe de analistas.

\section{Arquitetura \textit{Medallion} e Fragmentação de Ingestão}

Para organizar o \textit{pipeline} de maneira escalável, de fácil manutenção e rastreável, optou-se pela utilização da Arquitetura \textit{Medallion}. Trata-se de um padrão de projeto de organização de dados que divide os arquivos e processamentos em camadas lógicas, com o objetivo de melhorar progressivamente a estrutura e a qualidade das informações a cada nova etapa \cite{databricks_medallion, gujjala_medallion}.

Nesta arquitetura, a transição e o refinamento dos dados ocorrem na seguinte divisão:

\begin{itemize}
    \item \textbf{Camada \textit{Bronze} (L0):} Consiste na camada inicial de armazenamento, atuando como o ambiente de ingestão bruta (\textit{raw data}). Para lidar com a complexidade da extração de dados não estruturados (como os arquivos PDF abordados neste projeto), esta camada foi arquiteturalmente fragmentada em duas subzonas: a \textbf{L0 Original} e a \textbf{L0 Revisada}. 
    
    A \textbf{L0 Original} opera como uma \textit{landing zone} (zona de pouso) estrita, onde os arquivos são mantidos em seu estado original e imutável, garantindo o registro histórico exato e a possibilidade de reprocessamento seguro em caso de falhas \cite{gorelik19}. Já a \textbf{L0 Revisada} recebe esses dados em uma etapa subsequente. A literatura técnica estabelece que a camada \textit{Bronze} pode comportar um algoritmo de \textit{parsing} mínimo focado estritamente em garantir a consistência do formato de armazenamento, sem a aplicação de regras de negócio \cite{gujjala_medallion}. Assim, a L0 Revisada converte os dados do PDF para um formato semi-estruturado bruto (como tabelas extraídas), facilitando a transição sem aplicar os filtros de limpeza estrutural que são exclusivos da próxima camada.
    
    \item \textbf{Camada \textit{Silver} (L1):} Trata-se da camada intermediária. Nela, os dados tabulares provenientes da L0 Revisada são submetidos a rigorosos processos de limpeza, padronização, deduplicação e, principalmente, validação contra contratos predefinidos. É o estágio onde a informação deixa de ser um dado bruto e passa a ser um dado confiável e estruturado \cite{gujjala_medallion}.
    
    \item \textbf{Camada \textit{Gold} (L2):} É a camada final de processamento do repositório. Os dados, já limpos e validados, são agregados e transformados conforme as regras abrangentes da organização, ficando otimizados e prontos para o consumo analítico \cite{databricks_medallion}.
\end{itemize}

\section{Camada de Apresentação (\textit{Views}) e Versatilidade}

Embora a camada \textit{Gold} finalize o tratamento da informação, a entrega direta de tabelas físicas aos usuários analistas pode gerar engessamento. Para solucionar isso, as modernas arquiteturas de engenharia de dados adotam uma Camada de Apresentação (frequentemente materializada através de \textit{Views} lógicas) operando sobre a camada \textit{Gold} \cite{reis22}. 

Uma diretriz de projeto fundamental adotada neste trabalho foi a de propagar o máximo de campos e informações viáveis extraídas da fonte primária até a camada anterior à apresentação (L2). A premissa central de preservar a abrangência e a granularidade atômica dos dados garante uma extrema versatilidade ao \textit{pipeline} \cite{kimball13}. 

Se o cliente final (analista de mercado) alterar seus requisitos de negócio e solicitar a visualização de uma coluna extraída do PDF que não havia sido prevista nos relatórios iniciais, essa arquitetura evita retrabalho. Como o dado já foi extraído da L0 Revisada, validado na L1 e preservado até a camada L2, não há necessidade de reconstruir os \textit{scripts} de ingestão: basta editar os limites da \textit{View} para passar a expor a nova métrica. Esse desacoplamento entre armazenamento estruturado e apresentação final eleva a resiliência operacional da solução.

\section{Contratos de Dados e \textit{Schemas}}

Um dos principais desafios na ingestão de arquivos externos é a mudança inesperada na estrutura original dos dados, como o reposicionamento de valores em um PDF. Para evitar que essas alterações corrompam silenciosamente os relatórios finais da empresa, utilizam-se os Contratos de Dados (\textit{Data Contracts}), que são acordos técnicos estritos definindo o formato, a estrutura e a tipagem esperada das informações na transição entre as camadas do \textit{pipeline} \cite{kleppmann17, armbrust20}.

A principal ferramenta lógica para impor o cumprimento desses contratos é o \textit{Schema} (Esquema). O \textit{schema} especifica claramente os nomes das colunas, os tipos de dados exigidos e as restrições de obrigatoriedade \cite{liao20}. 

Ao se aplicar a imposição de esquemas (\textit{schema enforcement}) na transição entre a L0 Revisada e a camada \textit{Silver}, o sistema atua como uma barreira de segurança: registros que não cumprem os requisitos estruturais preestabelecidos são bloqueados ou sinalizados para correção. Essa prática é essencial para lidar com a evolução estrutural dos dados, protegendo o sistema contra quebras e garantindo que apenas informações íntegras fluam para as \textit{Views} de apresentação \cite{armbrust20}.